\section{The support cutoff and the two distinguished prime fibres}

% BEGIN H87P1
\paragraph{H87P1. The third period-height entry holds on the original family.}
Take any original parameter family $\mathcal F$, $x>\mathcal F.x_1$
and $i\in\mathcal F.I_x$. Use its actual arithmetic member $D$ and
size datum $L$. For a finite place $w$ write $p$ for its residue
characteristic. At every lift $j$, we have
\[
 p>X''(n)\quad\Longrightarrow\quad \mathrm{gain}_{j,w}=0.
\]
Outside the stored finite support this is the defining support
condition. Inside the support, primes in the signal or inert class use
the member's filled gain formula, whose first case sets the gain to
zero whenever $p>X''(n)$. For the remaining prime classes the
defining local condition sets every gain to zero. This argument does
not use a restriction on $u$ or an exception for any prime.

The error has a different support. For every word $u$ and lift $j$,
the defining evaluation of $E_{u,j,w}$ first tests $p=T_u$.
On the complementary branch it adds $v_p(D_{\Gamma,u})$ to a
case-defined correction. Its first two cases occur at $g_0,g_1$.
The ramified case requires $p=d$. The two remaining nonzero
corrections require both membership in the signal or inert class
and $p\le X''(n)$. Thus
\[
 p>X''(n),\quad p\notin\{T_u,g_0,g_1,d\},\quad
        p\nmid D_{\Gamma,u}
 \quad\Longrightarrow\quad E_{u,j,w}=0 .
\]
These evaluations apply to the stored $E$ because the family's
error-definition equality holds for that member and every lift.
Together the implications give the full third entry
$L.\mathrm{ZeroAbove}$ of $\mathrm{LinkDIV}_{\mathrm{at}}$.
The threshold $x_0=\mathcal F.x_1+1$ puts both implications in
the required eventual form, simultaneously for all indices. This
also applies on the design class of X87S4.
% END H87P1

% BEGIN H87P2
\paragraph{H87P2. Exact contribution above each distinguished prime.}
On the same member put $K=D.\mathrm{what}_2.LL$,
$e_w=v_w(p)$, $\ell_w=\log Nw/[K:\mathbb Q]$ and
$\mathrm{sz}(w,t)=2\ell_we_wt$.
Write $N=N_m$, $f=D.R.f$, $m=D.R.m$, $\eta=L.\mathrm{etaSz}$,
and $r=c.g\, f(N^2+m^2\eta)$ for the member's fixed error constant
$c.g$ from its family. Let $d_j$ be its Fock degree and put
$\nu_{u,j,\tau}=\mathrm{nuAbs}(D,u,j,\tau)$, the sum of its
stored row exponents. For each $g\in\{g_0,g_1\}$ the following
holds also at lifts whose degree difference vanishes.
\[
 \sum_{\substack{w\in L.\mathrm{supp}\\w\mid g}}
       \mathrm{sz}(w,E_{u,j,w})
  =\left(\sum_\tau\nu_{u,j,\tau}-d_j+
                     2v_g(D_{\Gamma,u})+2r\right)\log g .
\]

To see the normalization, for any $y\in K$ with $y\bar y=g$,
conjugation permutes the places above $g$ and preserves $\ell_w$.
The sums of $\ell_w v_w(y)$ and $\ell_w v_w(\bar y)$ over that
fibre agree. Their sum equals $\sum_{w\mid g}\ell_we_w=\log g$,
by $\sum_{w\mid g}e_wf_w=[K:\mathbb Q]$, where $f_w$ is the
residue degree. For every $\tau$ take $y$ to be the image under
$\tau$ of the conjugate of the chosen prime generator at $g$.
Its norm identity follows from $\pi_g\bar\pi_g=g$ and equivariance
of the coefficient embedding for conjugation. Hence
$2\sum_{w\mid g}\ell_wv_w(y)=\log g$ for each $\tau$.

Since $g\ne T_u$, evaluation on the fibre is
\[
 E_{u,j,w}=v_g(D_{\Gamma,u})-d_j/2+r+
              e_w^{-1}\sum_\tau\nu_{u,j,\tau}v_w(y_\tau).
\]
Multiply by $2\ell_we_w$ and sum. The norm equality evaluates
each last summand. Places outside the support add zero by the
support axiom for $E$. This proves the asserted identity with
the original weights and both factors of two.
% END H87P2

% BEGIN H87P3
\paragraph{H87P3. Reduction of the eighth entry to the numerical estimates.}
Let $A=\sum_{t\in TT}a_t$ and $B_0=fmbx^\epsilon$.
On the late range of a family in the design class of X87S4, the
level-prime and stabilizer-size bounds $\mathrm{TuSize}$ and
$\mathrm{DGamSize}$ are components of class membership. For
$0<\epsilon<1/52$ the pinned numerical route supplies six uniform
estimates. We describe their exact use in the error sum so the
remaining numerical inputs can be read independently of the periods.

First the degree-count estimate bounds
$|\sum_\tau\nu_{u,j,\tau}-d_j-A|$ by
$\mathrm{errG}(L,C_{\rm deg},j)$. The distinguished-prime size
estimate bounds the product
\[
  \big(\mathrm{errG}(L,C_{\rm deg},j)+2v_g(D_{\Gamma,u})+2r\big)
         \log g \le C_g B_0 .
\]
Here $r\ge0$ by the prescribed positive error constants. H87P2 now
bounds the difference from $A\log g$ at each of the two primes by
$C_gB_0$. This step uses their separate weighted sums, including
the local valuations of the stabilizer integer.

Partition the remaining places into those above $T_u$ and those
away from $T_u$. On the first part the nonnegative local evaluation
is $E_{T_u}$ and the level-prime estimate gives the bound $C_TB_0$
for its weighted sum. On the second part, writing $E_{d,j}$ for
the ramified correction, the triangle inequality in the evaluation
formula of H87P1 gives
\[
 |E_{u,j,w}|\le v_p(D_{\Gamma,u})
 +{\bf1}_{p=d}|E_{d,j}|
 +{\bf1}_{\substack{p\ \text{signal or inert}\\p\le X''(n)}}
                      |E_{CI,j}(p)|.
\]
Here $E_{CI,j}(p)$ equals the conjugation-prime correction on the
inert class with $h_P=0$ and the other inert-prime correction
otherwise. Each right-hand side depends only on $p$. As
$\sum_{w\mid p}e_w\ell_w=\log p$, a sum over any part of a
prime fibre is bounded by the full fibre for these nonnegative
majorants. Prime factorization bounds the first resulting sum by
$2\log D_{\Gamma,u}$. The other two sums are bounded by
$2\log d\,|E_{d,j}|$ and the sum of
$2\log p\,|E_{CI,j}(p)|$ over the indicated primes.
The stabilizer, ramified-prime and inert-prime estimates bound
these three quantities by $C_\Gamma B_0,C_d B_0,C_I B_0$.
Consequently the absolute sum away from $g_0,g_1$ is bounded
by $(C_T+C_\Gamma+C_d+C_I)B_0$. The primes $g_0,g_1$ are
distinct, so $A\log g_0+A\log g_1=A\log(g_0g_1)$. Adding the
three fibre estimates gives the original eighth entry with
\[
                       C=2C_g+C_T+C_\Gamma+C_d+C_I .
\]
Take the maximum of the six numerical thresholds and
$\mathcal F.x_1+1$. All words, lifts, and members then use this
same constant and threshold for the fixed family.

The numerical proof in the pinned source chooses $C_{\rm deg}=40$,
$C_T=\max(12C_{Tu},0)$, $C_d=4$, and
$C_I=18(1/\epsilon+1)+848$ by splitting the two
inert-prime corrections. The proof above supplies the exact
summation implication from the six estimates. Their detailed
numerical derivations, including the fine degree count and the two
inert-prime bounds, remain inputs to this manuscript passage.
% END H87P3
